% 天体（综述）
% license CCBYSA3
% type Wiki

本文根据 CC-BY-SA 协议转载翻译自维基百科\href{https://en.wikipedia.org/wiki/Astronomical_object}{相关文章}。

一个天文对象、天体、恒星对象或天界对象，是指在可观测宇宙中自然存在的物理实体、关联或结构。在天文学中，“对象”（object）和“天体”（body）这两个术语常被交替使用。然而，“天文天体”“天体”或“天界天体”特指单一的、紧密结合的、连续的物理体；而“天文对象”或“天体对象”则可以指更为复杂、结合程度较弱的结构，其中可能包含多个天体，甚至包含带有次级结构的其他对象。

天文对象的例子包括行星系统、星团、星云和星系；而小行星、卫星（moons）、行星和恒星则属于天文天体。一颗彗星可以同时被视为天体和对象：当指其由冰与尘埃组成的冻结核时，它是一个天体；而当描述其整体，包括弥散的彗发与彗尾时，它则属于天文对象。
